% Figure 12 — parity is a trajectory comparison, not a two-box architecture.
\begin{figure}[H]
\centering
\resizebox{0.89\textwidth}{!}{%
\begin{tikzpicture}[pd figure,x=1cm,y=1cm]
  \foreach \x/\op in {2.0/{open},4.25/{pragma},6.5/{write},8.75/{crash},11.0/{reopen},13.0/{read}} {
    \draw[pd guide] (\x,1.0)--(\x,5.95);
    \node[pd axis label,anchor=north] at (\x,.78) {\op};
  }
  \node[pd panel title,anchor=east] at (1.35,4.85) {test};
  \node[pd panel title,anchor=east] at (1.35,2.15) {deployment};
  \draw[pd focus rule,line width=1.35pt]
    (2.0,4.85)--(4.25,4.85)--(6.5,5.15)--(8.75,4.55)--(11.0,4.85)--(13.0,4.85);
  \draw[pd rule,line width=1.35pt]
    (2.0,2.15)--(4.25,2.15)--(6.5,2.45)--(8.75,1.85)--(11.0,2.15)--(13.0,2.15);
  \foreach \x/\y in {2.0/4.85,4.25/4.85,6.5/5.15,8.75/4.55,11.0/4.85,13.0/4.85}
    \node[pd focus datum] at (\x,\y) {};
  \foreach \x/\y in {2.0/2.15,4.25/2.15,6.5/2.45,8.75/1.85,11.0/2.15,13.0/2.15}
    \node[pd datum] at (\x,\y) {};
  \draw[pd caution arrow] (6.5,4.55)--(6.5,2.72);
  \fill[pd hatch] (5.85,2.62) rectangle (7.15,4.42);
  \node[pd direct label,anchor=west,text=hhamber] at (7.35,3.52)
    {binding semantics may diverge};
  \draw[pd focus rule] (12.7,4.42)--(13.3,4.42);
  \draw[pd focus rule] (12.7,2.58)--(13.3,2.58);
  \draw[pd focus arrow] (13.0,4.32)--(13.0,2.68);
  \node[pd direct label,anchor=west,text=hhteal] at (13.35,3.52)
    {compare observable state};
  \node[pd note,anchor=west] at (1.75,6.25)
    {same operation sequence; parity is earned only by matching outcomes};
\end{tikzpicture}%
}
\caption{Runtime parity as a trajectory comparison. The same operation sequence
is applied to the test binding and to the compiled deployment daemon; the harness
must compare observable state after reopen, not merely confirm that both programs
build. The hatched gap is the unproved surface where pragma, locking, or checkpoint
semantics can diverge. A shim aligns interfaces, but only booting the deployment
runtime and differential testing can collapse that gap into a verified invariant.}
\label{fig:swk-dual-runtime}
\end{figure}
